\subsection{Spektrum kvartického oscilátoru}
\label{sec:QuarticOscillator}
\begin{enumerate}
    \item
        V jednorozměrném harmonickém oscilátoru popsaném Hamiltoniánem
        \begin{equation}
            \operator{H}_{0}=\frac{1}{2M}\operator{p}^{2}+\frac{1}{2}M\Omega^{2}\operator{x}^{2}
        \end{equation}
        vyjádřete maticové elementy
        \begin{subequations}
            \begin{align}
                &\matrixelement{m}{\operator{x}^{2}}{n},\\
                &\matrixelement{m}{\operator{x}^{4}}{n},
                \end{align}                    
        \end{subequations}
        kde $\ket{m}$ a $\ket{n}$ jsou dva libovolné vlastní vektory Hamiltoniánu $\operator{H}_{0}$.
        
    \item
        Uvažujte Hamiltonián
        \begin{equation}
            \operator{H}=\frac{1}{2M}\operator{p}^{2}+\frac{1}{2}M\Omega^{2}\operator{x}^{2}+b\operator{x}^{4}
                =\operator{H}_{0}+b\operator{x}^{4},
        \end{equation}
        kde $b$ je reálný kladný parametr.
        
        \begin{itemize}
        \item
            Napočítejte maticové elementy $H_{mn}=\matrixelement{m}{\operator{H}}{n}$ pro $M=\Omega=\hbar=b=1$ a numerickou diagonalizací matice $\matrix{H}=\left(H_{mn}\right)$ 
            (například pomocí programů Mathematica, Maple, Matlab, GNU Octave, Python, Julia, knihoven LAPACK, atd.) učete energie základního stavu $E_{0}$ a prvních tří vzbuzených stavů $E_{1,2,3}$ Hamiltoniánu $\operator{H}$.
            Uvažujte pouze konečný počet $N$ elementů báze, tj. $m,n=0,1,\dotsc,N-1$.

        \item 
            Zakreslete graf závislosti $E_{j}(N)$, $j=0,1,2,3$ pro $N\leq50$.

        \item
            Na základě výsledků z předchozího bodu odhadněte, jaká musí být nejmenší velikost matice $N$, aby čtyři nejnižší energie byly určeny s přesností na pět platných cifer.
        \end{itemize}
        
    \item
        Uvažujte Hamiltonián
        \begin{equation}
            \label{eq:HamiltonianQuarticOscillator}
            \operator{H}'=\frac{1}{2M}\operator{p}^{2}-\frac{1}{2}a\operator{x}^{2}+\frac{1}{8}b\operator{x}^{4}\,,
        \end{equation}
        kde $a>0$, $b>0$.
        
        \begin{itemize}
        \item 
            Načtněte klasický potenciál
            \begin{equation}
                V(x)=-\frac{1}{2}ax^{2}+bx^{4}
            \end{equation}
            pro $a=b=1$ a nalezněte jeho stacionární body.
        \item 
            Vypočítejte numericky první čtyři energetické hladiny pro $\hbar=0.02,M=a=b=1$.
            Velikost matice ať je taková, aby tyto hladiny byly určeny s přesností alespoň na pět platných cifer.
        \item 
            Proč jsou dvojice hladin $(E_{0}, E_{1})$ a $(E_{2}, E_{3})$ v téměř degenerovaných dubletech?        
        \item 
            Opakujte výpočet pro jiné hodnoty $\hbar$, například $\hbar=0.01$ a $\hbar=0.1$, a diskutujte, jaký vliv má volba $\hbar$ na výsledné spektrum.
            Nezapomeňte pokaždé zkontrolovat konvergenci energetických hladin (obecně pro menší hodnoty $\hbar$ je potřeba více bázových stavů).
        \end{itemize}		
        
        \emph{Poznámka:} Dvoujámový systém popsaný Hamiltoniánem $\operator{H}'$ se požívá například k modelování amoniakového maseru, k modelování systémů ochlazených iontů, v teorii kvantové informace, při studiu kvantových fázových přechodů nebo v kvantové chemii.
            
    \end{enumerate}
